% Generated draft: requires mathematical and source review.
\documentclass[11pt]{article}
\usepackage[margin=25mm]{geometry}
\usepackage{fontspec}
\setmainfont{Latin Modern Roman}
\usepackage{amsmath,amssymb,mathtools,unicode-math}
\setmathfont{Latin Modern Math}
\usepackage{enumitem,needspace,xurl,hyperref}
\setlength{\parindent}{0pt}
\setlength{\parskip}{5pt}
\setlength{\emergencystretch}{3em}
\newcommand{\N}{\mathbb{N}}
\newcommand{\Z}{\mathbb{Z}}
\newcommand{\Q}{\mathbb{Q}}
\newcommand{\R}{\mathbb{R}}
\newcommand{\C}{\mathbb{C}}
\newcommand{\F}{\mathbb{F}}
\newcommand{\A}{\mathbb{A}}
\newcommand{\PP}{\mathbb{P}}
\newcommand{\E}{\mathbb{E}}
\newcommand{\eps}{\varepsilon}
\newcommand{\dd}{\,\mathrm{d}}
\begin{document}
\section*{UnsolvedMath ID: 20002907}
\textbf{Source problem code:} catalogue entry
\subsection{Definitions and mathematical statement}
The source considers the level-\(k\) fusion product of rectangular type A factors
\[
 B=B^{r_\ell,s_\ell}\otimes_k\cdots\otimes_k B^{r_1,s_1}
\]
and asks for a two-variable \((q,t)\) generalization of its grading. It reports type A combinatorial/fermionic fusion formulas through rigged configurations (Schilling et al.) and conjectural other-type formulas (Hatayama et al.). The target multiplicity, energy normalization and meaning of a second grading are not fixed; a positive generating polynomial alone is weaker than a module bigrading.

The aid uses type \(A_n^{(1)}\) rectangular crystals, \(1\le r_j\le n\), \(s_j\ge1\), and fixes a dominant classical weight \(\lambda\). Let \(L_i^{(a)}\) count factors \(B^{a,i}\). A configuration consists of partitions \(\nu^{(a)}\), with \(m_i^{(a)}\) parts of size \(i\), satisfying
\[
 \sum_{a,i}i m_i^{(a)}\alpha_a
       =\sum_{a,i}iL_i^{(a)}\omega_a-\lambda.
\]
For the type A Cartan matrix \(C\), define vacancy numbers and configuration cocharge by
\[
 p_i^{(a)}=\sum_j\min(i,j)L_j^{(a)}
      -\sum_b C_{ab}\sum_j\min(i,j)m_j^{(b)},
\]
\[
 \operatorname{cc}(\nu)=\tfrac12\sum_{a,b,i,j}
          C_{ab}\min(i,j)m_i^{(a)}m_j^{(b)}.
\]
For an admissible configuration, each rigging \(J_i^{(a)}\) is a partition fitting an \(m_i^{(a)}\times p_i^{(a)}\) rectangle. Put \(R(\nu,J)=\sum_{a,i}|J_i^{(a)}|\).

Fix the aid's energy-compatible rigged-configuration bijection \(\Psi\), obtained from the chosen standard bijection by the appropriate corigging complement, and normalize the path energy by
\[
 D(b)=\operatorname{cc}(\nu)+R(\nu,J),\qquad (\nu,J)=\Psi(b).
\]
This convention must be matched to the one-variable fusion polynomial; a change of additive energy normalization changes the prefactor.

At stable level \(k\ge\sum_j s_j\), the aid reports the positive refinement
\[
 F_{B,\lambda}(q,t)=\sum_{b\text{ highest weight }\lambda}
                     q^{D(b)}t^{R(\Psi(b))}
 =\sum_{\nu\text{ admissible}}q^{\operatorname{cc}(\nu)}
       \prod_{a,i}
       \begin{bmatrix}p_i^{(a)}+m_i^{(a)}\\m_i^{(a)}\end{bmatrix}_{qt}.
\]
The Gaussian binomial at \(u\) counts partitions in the indicated rectangle by \(u^{|J|}\). Thus the expression has nonnegative coefficients. In the reported normalization, \(t=1\) gives the usual one-dimensional \(X=M\) polynomial, and \(q=t=1\) gives the corresponding fusion coefficient. The second degree is invariant under changing factor order through the combinatorial \(R\)-matrix and the compatible rigged-configuration identification, not through arbitrary unrelated bijections.

For a fixed configuration put \(A_\nu=\sum_{a,i}m_i^{(a)}p_i^{(a)}\). Complementing each rectangle sends \(R\) to \(A_\nu-R\). If \(P_\nu(u)\) is the product of Gaussian binomials, the exact fiber reciprocity is
\[
 u^{A_\nu}P_\nu(u^{-1})=P_\nu(u).
\]
Equivalently the fiber \(q^{\operatorname{cc}(\nu)}P_\nu(qt)\) is reciprocal with factor \(q^{2\operatorname{cc}(\nu)+A_\nu}t^{A_\nu}\) under \((q,t)\mapsto(q^{-1},t^{-1})\). The factor depends on \(\nu\); no global single-degree reciprocity is asserted.

For arbitrary fusion level, let \(\mathrm{RC}_k\) be the established level-restricted set in the standard bijection convention, and let \(\theta\) complement riggings. The aid's energy-side set is \(\theta(\mathrm{RC}_k)\), not automatically \(\mathrm{RC}_k\) itself. Summing \(q^{\operatorname{cc}(\nu)+R}t^R\) over this image gives its finite-level positive candidate. Level restrictions can couple riggings, so the stable Gaussian product and complement reciprocity do not automatically persist. The summary does not supply the exact finite-level inequalities, a closed product formula, a representation-theoretic bigrading, or identification with a canonical Macdonald \((q,t)\) theory. Those remain additional parts of the original request.
\subsection{Short English statement}
The aid proposes total energy-side rigging size as a second statistic. At stable fusion level this gives a positive Gaussian-binomial formula with the correct one-variable and ungraded specializations. Finite-level restrictions require the complemented restricted set and need not factor. The construction is a combinatorial refinement, not a supplied module bigrading or canonical Macdonald deformation.
\subsection{Sources}
\begin{itemize}
\item Local UnsolvedMath registry, permanent ID 20002907, catalogue entry; canonical \texttt{aim-special-functions-notes.json}, notes[11], reconstructed-unverified statement status.
\item AIM \emph{Generalized Kostka polynomials}, Problem 12, \url{https://aimath.org/WWN/kostka/prob2.pdf}. No individual proposer is supplied.
\item Source attributions: Schilling et al. for type A rigged-configuration interpretations and fermionic fusion formulas; Hatayama et al. for conjectural other-type formulas. No full bibliographic entries are supplied, and the source's conjectural wording is a historical report rather than an independently updated status.
\item Machine-generated research attempt 1, partial-result/partially-solved labels, reported rigor proved and low novelty confidence. Its second statistic and stable-level formula are not peer reviewed and require independent verification.
\end{itemize}
\end{document}
